\SourceLeaf{61}
\begin{center}
\textit{FIFTH PROOF OF THE FUNDAMENTAL THEOREM}\par
\textit{IN THE THEORY OF QUADRATIC RESIDUES.}
\end{center}

\Article{1}

We have already stated in the introduction that the fifth and third proofs start from the same lemma. For convenience, we have thought it appropriate to repeat it here in notation adapted to the present investigation.

\textsc{Lemma.}\quad \textit{Let $m$ be a prime number (positive and odd), and $M$ an integer not divisible by $m$. Take the least positive residues of the numbers}

\[
\begin{aligned}
M,\ 2M,\ 3M,\ 4M\ \mathinner{.\,.\,.\,.\,.\,.}\ \tfrac12(m-1)M
\end{aligned}
\]

\textit{modulo $m$. Some will be less than $\tfrac12m$, and others greater; let the number of the latter be $=n$. Then $M$ is a quadratic residue or a nonresidue of $m$, according as $n$ is even or odd.}

\textsc{Proof.}\quad Let those of these residues which are less than $\tfrac12m$ be $a,b,c,d$ etc., and the remaining ones, greater than $\tfrac12m$, be $a',\ b',\ c',\ d'$ etc. The complements of the latter to $m$, namely $m-a',\ m-b',\ m-c',\ m-d'$ etc., are evidently all less than $\tfrac12m$ and distinct both from one another and from the residues $a,\ b,\ c,\ d$ etc. Thus, taken together with these, they are precisely all the numbers $1,\ 2,\ 3,\ 4\dotsfour\tfrac12(m-1)$, though in a different order. Putting the product

\[
\begin{aligned}
1\pdot2\pdot3\pdot4\dotsfive\tfrac12(m-1)=P
\end{aligned}
\]

we therefore have

\[
\begin{aligned}
P=abcd\dotsfour\times(m-a')(m-b')(m-c')(m-d')\dotsfour
\end{aligned}
\]

and hence

\[
\begin{aligned}
(-1)^nP=abcd\dotsfour\times(a'-m)(b'-m)(c'-m)(d'-m)\dotsfour
\end{aligned}
\]

Furthermore, modulo $m$, we have

\[
\begin{aligned}
PM^{\frac12(m-1)}&\equiv abcd\dotstwo\times a'b'c'd'\dotstwo\\&\equiv abcd\dotstwo\times(a'-m)(b'-m)(c'-m)(d'-m)\dotstwo
\end{aligned}
\]

and hence

\[
\begin{aligned}
PM^{\frac12(m-1)}\equiv P(-1)^n
\end{aligned}
\]

It follows that $M^{\frac12(m-1)}\equiv\pm1$, taking the upper or lower sign according as $n$ is even or odd. With the aid of the theorem proved in Article 106 of the \textit{Disquisitiones Arithmeticae}, the lemma follows immediately.

\SourceLeaf{62}
\Article{2}

\textsc{Theorem.}\quad \textit{Let $m,M$ be positive odd integers coprime to one another, and let $n$ be the number of those among the least positive residues of the numbers}

\[
\begin{aligned}
M,\ 2M,\ 3M\dotsfive\tfrac12(m-1)M
\end{aligned}
\]

\textit{modulo $m$ which are greater than $\frac12m$. Likewise, let $N$ be the number of those among the least positive residues of the numbers}

\[
\begin{aligned}
m,\ 2m,\ 3m\dotsfive\tfrac12(M-1)m
\end{aligned}
\]

\textit{modulo $M$ which are greater than $\frac12M$.\quad Then the three numbers $n,N$, $\frac14(m-1)(M-1)$ are either all even, or one is even and the other two odd.}

\textsc{Proof.}\quad Let us denote

{\small
\begin{center}\setlength{\tabcolsep}{3pt}
\begin{tabular}{@{}lll@{}}
by $f$ & the collection of numbers & $1,\ 2,\ 3\dotsfive\frac12(m-1)\mathord{.}$ \\
by $f'$ & the collection of numbers & $m-1,\ m-2,\ m-3\dotsfive\frac12(m+1)$ \\
by $F$ & the collection of numbers & $1,\ 2,\ 3\dotsfive\frac12(M-1)$ \\
by $F'$ & the collection of numbers & $M-1,\ M-2,\ M-3\dotsfive\frac12(M+1)$
\end{tabular}\end{center}
}

Thus $n$ indicates how many numbers $Mf$ have their least positive residues modulo $m$ in the collection $f'$. Likewise, $N$ indicates how many numbers $mF$ have their least positive residues modulo $M$ in the collection $F'$. Finally, let

{\small
\begin{center}\setlength{\tabcolsep}{3pt}
\begin{tabular}{@{}lll@{}}
$\varphi$ & denote the collection of numbers & $1,\ 2,\ 3\dotsfive\frac12(mM-1)$ \\
$\varphi'$ & denote the collection of numbers & $mM-1,\ mM-2,\ mM-3\dotsfive\frac12(mM+1)$
\end{tabular}\end{center}
}

Every integer not divisible by $m$ must be congruent modulo $m$ either to a residue from $f$ or to one from $f'$. Likewise, every integer not divisible by $M$ is congruent modulo $M$ either to a residue from $F$ or to one from $F'$. Thus all the numbers in $\varphi$, among which evidently none is divisible by both $m$ and $M$, can be divided into eight classes as follows.

I.\quad The first class contains the numbers congruent modulo $m$ to a number from $f$, and modulo $M$ to a number from $F$.\quad We shall denote their number by $\alpha$.

\SourceLeaf{63}
II.\quad The numbers congruent modulo $m,M$, respectively, to numbers from $f,F'$; we put their number $=\beta$.

III.\quad The numbers congruent modulo $m,M$, respectively, to numbers from $f',F$; we put their number $=\gamma$.

IV.\quad The numbers congruent modulo $m,M$, respectively, to numbers from $f',F'$; let their number be $=\delta$.

V.\quad The numbers divisible by $m$, but congruent modulo $M$ to residues from $F$.

VI.\quad The numbers divisible by $m$, but congruent modulo $M$ to residues from $F'$.

VII.\quad The numbers divisible by $M$, but congruent modulo $m$ to residues from $f$.

VIII.\quad The numbers divisible by $M$, but congruent modulo $m$ to residues from $f'$.

Evidently classes V and VI together contain all the numbers $mF$. The number contained in VI is $=N$, so the number contained in V is $\frac12(M-1)-N$.\quad Likewise, classes VII and VIII together contain all the numbers $Mf$; class VIII contains $n$ numbers, and class VII contains $\frac12(m-1)-n$.

In just the same way, all the numbers in $\varphi'$ are divided into eight classes IX\,.\,.\,XVI. If we keep the same order in doing this, it is easy to see that the numbers in the classes

\begin{center}IX, X, XI, XII, XIII, XIV, XV, XVI\end{center}

are, respectively, the complements to $mM$ of the numbers in the classes

\begin{center}IV, III, II, I, VI, V, VIII, VII\end{center}

so that class IX contains $\delta$ numbers, class X contains $\gamma$, and so on.\quad It is now clear that, if all the numbers of the first class are combined with all those of the ninth, we obtain all the numbers below $mM$ which are congruent modulo $m$ to a number from $f$, and modulo $M$ to a number from $F$. It is easy to see that their number equals the number of all combinations of each individual $f$ with each individual $F$.\quad We therefore have

\[
\begin{aligned}
\alpha+\delta=\tfrac14(m-1)(M-1)
\end{aligned}
\]

\SourceLeaf{64}
and by a similar argument we also have

\[
\begin{aligned}
\beta+\gamma=\tfrac14(m-1)(M-1)
\end{aligned}
\]

Combining all the numbers of classes II, IV and VI evidently gives all the numbers below $\frac12mM$ which are congruent to a residue from $F'$ modulo $M$.\quad These same numbers can also be written as follows:

\[
\begin{aligned}
F',\ M+F',\ 2M+F',\ 3M+F'\dotsfour\tfrac12(m-3)M+F'
\end{aligned}
\]

Their total number is therefore $=\frac14(m-1)(M-1)$, or we have

\[
\begin{aligned}
\beta+\delta+N=\tfrac14(m-1)(M-1)
\end{aligned}
\]

Likewise, combining all of classes III, IV and VIII gives

\[
\begin{aligned}
\gamma+\delta+n=\tfrac14(m-1)(M-1)
\end{aligned}
\]

\Needspace{9\baselineskip}
These four equations give the following:

\[
\begin{aligned}
2\alpha&=\tfrac14(m-1)(M-1)+n+N\\
2\beta&=\tfrac14(m-1)(M-1)+n-N\\
2\gamma&=\tfrac14(m-1)(M-1)-n+N\\
2\delta&=\tfrac14(m-1)(M-1)-n-N
\end{aligned}
\]

each of which establishes the theorem.

\Article{3}

If we now suppose that $m$ and $M$ are prime numbers, the fundamental theorem follows immediately by combining the preceding theorem with the lemma of Article 1.\quad For it is clear that

I.\quad whenever the numbers $m,M$ both have the form $4k+1$, or just one of them does, the number $\frac14(m-1)(M-1)$ is even. Thus $n$ and $N$ are either both even or both odd, and consequently either each of $m$ and $M$ is a quadratic residue of the other, or each is a quadratic nonresidue of the other.

II.\quad Whenever both numbers $m,M$ are of the form $4k+3$, however, \mbox{$\frac14(m-1)(M-1)$} is odd. Hence one of the numbers $n,N$ is even and the other odd, and consequently one of the numbers $m,M$ is a quadratic residue of the other, and the other a quadratic nonresidue of the first. Q. E. D.
